\documentclass[oneside, a4paper]{article}
\usepackage{kotex}
\begin{document}
\section[시작]{\LaTeX 시작하기}  
\begin{center}
안녕, \LaTeX!
\end{center}
한 번 배워서 오래 써먹자. \par  % 문단 구분 명령
이것은 첫 번째 문단의 첫 번째 문장입니다. 두 번째 문장입니다.

두 번째 문단의 첫 번째 문장입니다. 공백을 많이 쓴다고
공백이 그만큼 생기지 않습니다. 나쁜 습관입니다.

특수문자 \# \$ \% \& \_ \{ \} \string~ \string^ \textbackslash 를
입력할 때는 앞에 백슬래시를 더해주면 된다. 
\%는 주석(comment)를 나타내는 특수문자. \LaTeX이 컴파일
할 때 무시해 버린다 

`\LaTeX 에서는 여는 따옴표와 닫는 따옴표를 구분한다.'

터미널에 texdoc 명령을 입력하면 설명서를 쉽게 열 수 있다.
대부분의 패키지나 클래스는 texdoc 패키지 이름과 같은 방식으로 열 수 있다. 
예컨대, `texdoc kotex'을 입력하면 `kotexdoc.pdf' 파일이 열린다. 
\end{document}